\textbf{Phần 1. Trắc nghiệm lựa chọn}
\setcounterpageref{ex}{0}

\begin{ex} %[2D8V40]
Kết quả của phép tính $5 \times 7$ là
\choice
{$30$}
{$40$}
{$32$}
{\True $35$}
\loigiai{
$5 \times 7 = 35$.
\textbf{Đáp án của câu hỏi là D}
}
\end{ex}

\begin{ex} %[2D8V40]
Mỗi bàn tay có $5$ ngón tay. Hỏi $6$ bàn tay có bao nhiêu ngón tay?
\choice
{$25$ ngón tay}
{\True $30$ ngón tay}
{$35$ ngón tay}
{$20$ ngón tay}
\loigiai{
$5 \times 6 = 30$ (ngón tay).
\textbf{Đáp án của câu hỏi là B}
}
\end{ex}

\begin{ex} %[2D8V40]
Tính giá trị của biểu thức: $5 \times 8 - 15$.
\choice
{$20$}
{$30$}
{\True $25$}
{$35$}
\loigiai{
$5 \times 8 = 40$; $40 - 15 = 25$.
\textbf{Đáp án của câu hỏi là C}
}
\end{ex}

\begin{ex} %[2D8V40]
Một tuần lễ em đi học $5$ ngày. Hỏi trong $4$ tuần lễ em đi học bao nhiêu ngày?
\choice
{\True $20$ ngày}
{$18$ ngày}
{$22$ ngày}
{$25$ ngày}
\loigiai{
$5 \times 4 = 20$ (ngày).
\textbf{Đáp án của câu hỏi là A}
}
\end{ex}

\begin{ex} %[2D8V40]
Các kết quả trong bảng nhân 5 có chữ số tận cùng là
\choice
{Chỉ có chữ số $0$}
{Chỉ có chữ số $5$}
{Các chữ số chẵn}
{\True $0\text{ hoặc } 5$}
\loigiai{
Mọi bội số của 5 đều có tận cùng là 0 hoặc 5.
\textbf{Đáp án của câu hỏi là D}
}
\end{ex}

\begin{ex} %[2D8V40]
Tìm $x$ biết: $x : 5 = 6$.
\choice
{$25$}
{\True $30$}
{$35$}
{$20$}
\loigiai{
$x = 6 \times 5 = 30$.
\textbf{Đáp án của câu hỏi là B}
}
\end{ex}

\begin{ex} %[2D8V40]
Có $5$ túi gạo, mỗi túi nặng $5\text{ kg}$. Hỏi cả $5$ túi nặng bao nhiêu ki-lô-gam?
\choice
{$20\text{ kg}$}
{$30\text{ kg}$}
{\True $25\text{ kg}$}
{$24\text{ kg}$}
\loigiai{
$5 \times 5 = 25\text{ (kg)}$.
\textbf{Đáp án của câu hỏi là C}
}
\end{ex}

\begin{ex} %[2D8V40]
Mỗi can chứa $5\text{ L}$ dầu. Hỏi $9$ can như thế chứa bao nhiêu lít dầu?
\choice
{\True $45\text{ L}$}
{$40\text{ L}$}
{$50\text{ L}$}
{$35\text{ L}$}
\loigiai{
$5 \times 9 = 45\text{ (L)}$.
\textbf{Đáp án của câu hỏi là A}
}
\end{ex}

\begin{ex} %[2D8V40]
Điền số thích hợp: $5 \times 4 + 5 = 5 \times \dots$
\choice
{$4$}
{$6$}
{$7$}
{\True $5$}
\loigiai{
$5 \times 4 + 5 = 5 \times 5$. Vậy điền số 5.
\textbf{Đáp án của câu hỏi là D}
}
\end{ex}

\begin{ex} %[2D8V40]
Một ngôi sao $5$ cánh có $5$ đỉnh. Hỏi $7$ ngôi sao như thế có bao nhiêu đỉnh?
\choice
{$30$ đỉnh}
{\True $35$ đỉnh}
{$40$ đỉnh}
{$25$ đỉnh}
\loigiai{
$5 \times 7 = 35$ (đỉnh).
\textbf{Đáp án của câu hỏi là B}
}
\end{ex}

\begin{ex} %[2D8V40]
Tính nhanh: $5 \times 3 + 5 \times 7$.
\choice
{$45$}
{$55$}
{\True $50$}
{$40$}
\loigiai{
$5 \times (3 + 7) = 5 \times 10 = 50$.
\textbf{Đáp án của câu hỏi là C}
}
\end{ex}

\begin{ex} %[2D8V40]
Bác An nuôi $8$ chuồng thỏ, mỗi chuồng có $5$ con. Bác đã bán đi $12$ con thỏ. Bác còn lại bao nhiêu con thỏ?
\choice
{\True $28$ con}
{$26$ con}
{$30$ con}
{$24$ con}
\loigiai{
Tổng số thỏ: $5 \times 8 = 40$ (con).\nCòn lại: $40 - 12 = 28$ (con).
\textbf{Đáp án của câu hỏi là A}
}
\end{ex}

\textbf{Phần 2. Trắc nghiệm đúng sai}
\setcounterpageref{ex}{0}

\begin{ex} %[2D8V40]
Xét tính đúng sai của các phép tính trong bảng nhân 5:
\choiceTF
{\True Tích của 5 với số chẵn luôn có tận cùng là chữ số 0.}
{\True Tích của 5 với số lẻ luôn có tận cùng là chữ số 5.}
{Phép tính 5 x 6 có kết quả bằng 35.}
{\True Hai tích liên tiếp trong bảng nhân 5 hơn kém nhau 5 đơn vị.}
\loigiai{
\begin{itemize}
    \item Ý a) Đúng: Ví dụ: $5 \times 2 = 10, 5 \times 4 = 20$. Đúng.
    \item Ý b) Đúng: Ví dụ: $5 \times 3 = 15, 5 \times 5 = 25$. Đúng.
    \item Ý c) Sai: $5 \times 6 = 30$ (không phải 35).
    \item Ý d) Đúng: Đúng.
\end{itemize}
\begin{itemize}
    \item a) \textbf{Đúng}.
    \item b) \textbf{Đúng}.
    \item c) \textbf{Sai}.
    \item d) \textbf{Đúng}.
\end{itemize}
}
\end{ex}

\begin{ex} %[2D8V40]
Có 6 gói bánh, mỗi gói có 5 chiếc bánh quy:
\choiceTF
{\True Có tất cả 30 chiếc bánh quy.}
{Nếu chia đều 30 chiếc bánh cho 4 bạn thì mỗi bạn được 8 chiếc.}
{\True Nếu ăn hết 5 chiếc bánh thì số bánh còn lại bằng 5 x 5.}
{Số bánh ban đầu là một số lẻ.}
\loigiai{
\begin{itemize}
    \item Ý a) Đúng: $5 \times 6 = 30$ chiếc.
    \item Ý b) Sai: $30$ không chia hết cho 4 ($30 : 4 = 7$ dư 2).
    \item Ý c) Đúng: $30 - 5 = 25 = 5 \times 5$. Đúng.
    \item Ý d) Sai: 30 là số chẵn.
\end{itemize}
\begin{itemize}
    \item a) \textbf{Đúng}.
    \item b) \textbf{Sai}.
    \item c) \textbf{Đúng}.
    \item d) \textbf{Sai}.
\end{itemize}
}
\end{ex}

\textbf{Phần 3. Trả lời ngắn}
\setcounterpageref{ex}{0}

\begin{ex} %[2D8V40]
Tính kết quả phép tính: 5 x 8.
\shortans{40}
\loigiai{
\textbf{Đáp số:} $40$\\
$5 \times 8 = 40$.
}
\end{ex}

\begin{ex} %[2D8V40]
Một can đựng 5 lít nước mắm. Hỏi 7 can như thế đựng bao nhiêu lít nước mắm?
\shortans{35}
\loigiai{
\textbf{Đáp số:} $35$\\
$5 \times 7 = 35\text{ L}$.
}
\end{ex}

\begin{ex} %[2D8V40]
Tính giá trị: 5 x 6 + 25.
\shortans{55}
\loigiai{
\textbf{Đáp số:} $55$\\
$5 \times 6 = 30$; $30 + 25 = 55$.
}
\end{ex}

\begin{ex} %[2D8V40]
Mỗi bông hoa có 5 cánh hoa. Hỏi một bó gồm 9 bông hoa có tất cả bao nhiêu cánh hoa?
\shortans{45}
\loigiai{
\textbf{Đáp số:} $45$\\
$5 \times 9 = 45$ cánh hoa.
}
\end{ex}

\begin{ex} %[2D8V40]
Tìm số tự nhiên y biết: y : 5 = 8.
\shortans{40}
\loigiai{
\textbf{Đáp số:} $40$\\
$y = 8 \times 5 = 40$.
}
\end{ex}

\begin{ex} %[2D8V40]
Mẹ mua 5 hộp bút chì màu, mỗi hộp có 6 chiếc. Mẹ cho em 8 chiếc. Hỏi mẹ còn lại bao nhiêu chiếc bút chì màu?
\shortans{22}
\loigiai{
\textbf{Đáp số:} $22$\\
$5 \times 6 = 30$; $30 - 8 = 22$ chiếc.
}
\end{ex}
